\documentclass[11pt,a4paper]{article}
% Common preamble for the hint sheets (\input from each file)
\usepackage[utf8]{inputenc}
\usepackage[T1]{fontenc}
\usepackage{lmodern}
\usepackage{amsmath,amssymb,amsthm,mathtools}
\usepackage[a4paper,margin=2.3cm]{geometry}
\usepackage{enumitem}
\usepackage{xcolor}
\usepackage[hidelinks]{hyperref}
\renewcommand{\leq}{\leqslant}
\renewcommand{\le}{\leqslant}
\renewcommand{\geq}{\geqslant}
\renewcommand{\ge}{\geqslant}
\setlength{\parindent}{0pt}
\setlength{\parskip}{0.35em}
% \hintstitle{Sheet title}{Class work / Homework}
\newcommand{\hintstitle}[2]{\begin{center}{\LARGE\bfseries Hints: #1\par}\vspace{0.5em}{\large #2\par}\end{center}
\noindent\textit{The hints to each problem are cumulative. Read one, try again, and open the next only if you are still stuck.}\par\medskip\hrule\medskip}
\newcommand{\problem}[1]{\par\bigskip{\large\bfseries Problem #1}\par\nopagebreak\smallskip}
\newcommand{\subpart}[1]{\par\smallskip\textbf{(#1)}\par\nopagebreak}
\newcommand{\hint}[1]{\par\smallskip\noindent\textcolor{black!70}{\textbf{Hint #1.}}\enspace}
\newcommand{\note}[1]{\par\smallskip\noindent\textit{Note.} #1\par}
\providecommand{\lt}{<}
\providecommand{\gt}{>}
\providecommand{\tightlist}{\setlength{\itemsep}{0pt}\setlength{\parskip}{0pt}}
\emergencystretch=3em

\begin{document}
\hintstitle{Abstract algebra -- 4. Group actions, homomorphisms and automorphisms}{Homework}

\problem{1}
\hint{1} Two easy facts: $N_{UVW}=N_{VWU}=N_{WUV}$, since $xyz=e\iff yzx=e$; and $N_{UVG}=|U|\,|V|$, since $z$ is determined by $x,y$.
\hint{2} Expand $|A|\,|B|=N_{ABG}=N_{ABA}+N_{ABB}+N_{ABC}$ and $|B|\,|A|=N_{BAG}=N_{BAA}+N_{BAB}+N_{BAC}$. The first two terms match by cyclic invariance, so $N_{ABC}=N_{BAC}$, which is $N_{CBA}$ by cyclic invariance.

\problem{2}
\hint{1} Let $T=\{gh,gh^{-1},g^{-1}h,g^{-1}h^{-1}\}$. You must show that every element is a product of between $1$ and $|G|$ elements of $T$. Let $S$ be the set of all finite products of elements of $T$. In a finite group, a set closed under multiplication is a subgroup.
\hint{2} $S$ contains $(gh)(g^{-1}h)^{-1}=g^2$. Since $\operatorname{ord}(g)$ is odd, $g\in\langle g^2\rangle\subseteq S$, then $h=g^{-1}(gh)\in S$. So $S=G$.
\hint{3} For the bound, let $U_r$ be the set of products of at most $r$ elements of $T$, so $U_{r+1}=T\cup U_rT$. If $U_{r+1}=U_r$, the chain is stationary from then on. Before that it grows strictly, starting from $|U_1|\ge1$, so $U_{|G|}=S=G$.

\problem{3}
\hint{1} An endomorphism is determined by the images of a generating set. Choose generators greedily, $g_{i+1}\notin\langle g_1,\dots,g_i\rangle$. Each step raises the order of the generated subgroup by a factor that is at least a prime, so $G$ has a generating set with $k\le\Omega(n)$ elements, where $\Omega$ counts prime factors with multiplicity.
\hint{2} So $|\operatorname{End}G|\le n^{k}\le n^{\Omega(n)}$. Write $n=pm$; then $\Omega(n)=1+\Omega(m)\le 1+(m-1)=n/p$, because $\Omega(m)\le\log_2m\le m-1$.
\hint{3} Equality needs $\Omega(m)=m-1$, i.e.\ $m\in\{1,2\}$, and every choice of images of the generators must extend. So $n=p$ or $n=2p$ (with $p$ the largest prime). Check the candidates: $\mathbb Z/p$ has $p=n^{1}$ endomorphisms; for $n=4$, $(\mathbb Z/2)^2$ has $16=4^2$ and $\mathbb Z/4$ only $4$; for $n=2p$ with $p$ odd, $\mathbb Z/2p$ has $2p$ and $D_p$ has $p^2+1$, both less than $4p^2$. Equality: $\mathbb Z/p$ and $(\mathbb Z/2)^2$.

\problem{4}
\hint{1} Let $d=\gcd(m,n)=sm+tn$, $A=G(m)$, $B=G(n)$. Both are normal (generated by conjugation-invariant sets). Since $x^d=(x^m)^s(x^n)^t$, you get $G(d)=\langle A,B\rangle$.
\hint{2} For $a\in A$, $b\in B$ the commutator $[a,b]$ lies in $A\cap B$, by normality of each. Elements of $A\cap B$ commute with $A$ and with $B$, hence are central in $H=G(d)$. So $H$ has nilpotency class at most $2$, and the commutator is bimultiplicative on $H$.
\hint{3} Let $p=x^d$, $q=y^d$, and $m=dm'$, $n=dn'$. Then $p^{m'},q^{m'}\in A$ commute, so $[p,q]^{m'^2}=e$; likewise $[p,q]^{n'^2}=e$. Since $\gcd(m'^2,n'^2)=1$, $[p,q]=e$. So the generators of $H$ commute.

\problem{5}
\hint{1} An endomorphism of $\mathbb Z/n$ is $x\mapsto ax$ with $a=f(1)$, and it is bijective iff $\gcd(a,n)=1$. An additive map of $(\mathbb Z/p)^k$ is $\mathbb F_p$-linear. $|\mathrm{GL}_k(\mathbb F_p)|=\prod_{i=0}^{k-1}(p^k-p^i)$: choose the columns one by one outside the span of the previous ones.
\hint{2} If $|\operatorname{Aut}G|=1$, then $\operatorname{Inn}G\cong G/Z(G)$ is trivial, so $G$ is abelian. Then $x\mapsto x^{-1}$ is an automorphism, so $x^2=e$ for all $x$.
\hint{3} So $G\cong(\mathbb Z/2)^k$ with $\mathrm{GL}_k(\mathbb F_2)$ trivial, i.e.\ $k\le1$: $G$ is trivial or $\mathbb Z/2$.

\problem{6}
\hint{1} The rotation group of the cube has $24$ elements: the identity; $6$ rotations by $\pm90^\circ$ about face axes; $3$ by $180^\circ$ about face axes; $8$ by $\pm120^\circ$ about vertex axes; $6$ by $180^\circ$ about edge axes. For each type, count the cycles on faces and on vertices.
\hint{2} Faces: $\frac1{24}(k^6+3k^4+12k^3+8k^2)$. Vertices: $\frac1{24}(k^8+17k^4+6k^2)$. Check: for $k=2$ these give $10$ and $23$.
\hint{3} Integrality follows from Burnside: the expression counts orbits. A polynomial with non-integral coefficients can still be integer-valued, like $\binom k2$.

\problem{7}
\hint{1} Let $e$ be the identity of $(G,*)$. For any $a$, $a*e=a$, while $a\times e\perp a$. So either $a=0$ or $a\times e=0$. If $e\ne0$, every element is parallel to $e$ and all cross products vanish. So assume $e=0$.
\hint{2} Suppose $a\times b\ne0$, so $c=a*b=a\times b$. Compute $(a*a)*b=a*c=a\times c$ and $(a*a)*c=a*(a\times c)=a\times(a\times c)=-|a|^2c$. Neither is orthogonal to $b$, resp.\ $c$.
\hint{3} So $(a*a)\times b=0$ and $(a*a)\times c=0$, i.e.\ $a*a$ is parallel to both $b$ and $c$, hence $a*a=0=e$. Then $b=(a*a)*b=a\times c=a\times(a\times b)$, which forces $a\perp b$ and $b=-|a|^2b$: impossible.

\problem{8}
\hint{1} Let $F(g)$ be the number of fixed points and $n=|X|$. Burnside gives $\sum_gF(g)=|G|$. Applying it to the action on $X\times X$ (at least two orbits: the diagonal and its complement) gives $\sum_gF(g)^2\ge2|G|$.
\hint{2} Consider $\sum_g(F(g)-1)(F(g)-n)$. It equals $\sum F^2-|G|\ge|G|$. Each term with $1\le F(g)\le n$ is $\le0$, and each $g$ without fixed points contributes exactly $n$.
\hint{3} Hence $n\cdot\#\{g:F(g)=0\}\ge|G|$, so the proportion is at least $1/n$, in particular positive.

\problem{9}
\hint{1} $\varphi\circ i_g\circ\varphi^{-1}=i_{\varphi(g)}$.
\hint{2} Two transpositions commute iff they are disjoint (otherwise their product has order $3$). So $\varphi(1\,2),\dots,\varphi(1\,n)$ are pairwise non-commuting transpositions. For $n\ge4$, show they all share one common point $a$ (three pairwise intersecting transpositions either share a point or form a triangle $(xy),(yz),(xz)$, and a fourth one cannot meet all three sides of a triangle).
\hint{3} Write $\varphi(1\,i)=(a\,b_i)$ and define $\sigma(1)=a$, $\sigma(i)=b_i$. Then $\varphi$ and conjugation by $\sigma$ agree on the generators $(1\,i)$. ($n\le3$: check directly.)

\problem{10}
\hint{1} The action on the $n$ cosets gives $G\to S_n$. Its kernel $K$ is normal, and $G/K$ embeds in $S_n$, so $[G:K]$ divides $n!$.
\hint{2} If a simple group of order $60$ had a subgroup of index $k$ with $1<k<5$, the kernel would be a proper normal subgroup, hence trivial. Then $G$ would embed in $S_k$, and $60\mid k!$ is impossible for $k\le4$.

\end{document}
